\section{Data and Methodology}\label{sec:data}

\subsection{Sources and population}\label{sec:sources}

The data are a frozen extract of two PGE-SP tables.
\texttt{DEBITO} holds monthly snapshots of enrolled certificates.
\texttt{ARRECADACAO} holds GARE payments linked to the certificate.
The cross-section is the December 2025 snapshot.
Payments are those dated before 1~January 2026.
No 2026 debt and no 2026 payment enter the results.
The unit of the portfolio is the certificate.
The firm is the eight-digit root of the CNPJ.
No name, no individual tax number and no certificate identifier leaves the computation.
Cells with one to four firms are left blank in the detailed activity files that sit behind the paper.

The December 2025 extract of legal-entity \vat{} debt contains 2,325,277 certificates and \brl{}427.880 billion.
The project's published laboratory page reports 2,325,280 certificates and the same stock at the published precision.
The gap is three certificates and about \brl{}68,000.
We use the extract.
Inside it, declared debt is 2,257,189 certificates and \brl{}116.42 billion.
Assessed debt is 59,542 certificates and \brl{}311.40 billion, of which the Simples Nacional assessments are 818 certificates and \brl{}0.145 billion, or 0.05~percent of assessed stock.
Fecoep is 8,510 certificates and \brl{}0.048 billion.
The GIA amendment is 36 certificates and \brl{}0.009 billion.
Those two types stay out of every declared-versus-assessed comparison below.

We do not use changes between monthly snapshots.
Some months in the source archive are missing or partial, including December 2018 and October to December 2020, and the December 2024 snapshot is partial.
A panel of stocks would need a separate construction.
December 2025 in the extract is one row per certificate.
The results use that snapshot and the dates of payments.

Firm attributes come from the Receita Federal CNPJ registry already linked to these debtors, matched on the 14-character CNPJ of the establishment.
More than 99~percent of certificates in each segment match an establishment.
Activity is the principal CNAE of that establishment, placed on the official CNAE-Subclasses 2.3 hierarchy of section, division, group, class and subclass \citep{ibge2019cnae}.
The hierarchy file contains 1,331 subclasses.
Every debtor subclass is in that file except the placeholder code 8888888, which covers 808 certificates and is not an activity.
Section labels in the figures are short English renderings.
Division names in the table follow the official Portuguese titles, shortened in English only where the column must fit.
IBGE does not publish an official English list of the divisions.
Size, legal form, registry status, state and municipality are taken from the same register.
Opening date is the establishment's start of activity.

The register is observed after the debts and after the payments.
Status, activity, size and age describe the firm as the register stands.
They are not traits measured before the outcome, and a firm may be closed or unfit because of the debt.
We report them as descriptors and we do not give them a causal reading.

\subsection{Measures}\label{sec:measures}

Stock is the certificate's value net of legal fees in the December 2025 snapshot.
The payment rate is the share of certificates with any GARE payment dated before 1~January 2026.
Recovery is the sum of those payments divided by the December 2025 stock.
It is reais collected per real of stock still enrolled.
It is not an annual hazard, and it is not payments divided by the sum of past stocks.
A large assessed balance can show some payment and still a very small recovery.
Debt age is the number of years from the enrolment date to 31~December 2025.
Enrolment dates before 1960 are treated as missing.
In this snapshot the age comparison does not lose certificates to that rule.
Firm age is the number of years from the opening date to the same day, for establishment matches, with opening dates before 1900 treated as missing.
Where we report a firm-level age, it is the stock-weighted mean of the firm's certificates in that segment.

Concentration is computed on certificates with a positive value.
The Lorenz curve orders certificates from the smallest value to the largest and plots the cumulative share of stock against the cumulative share of certificates \citep{lorenz1905}.
The Gini coefficient is one minus twice the area under that curve.
The top 1~percent is the largest $\lceil 0.01 \times N \rceil$ positive-value certificates.
We use that round cut for two reasons, which Section~\ref{sec:conc} reports rather than assumes.
Statistically, the pooled Lorenz curve is flat over most of the distribution and rises in the extreme tail, so 1~percent is a visible break and not an interior point on a smooth curve.
Administratively, the resulting list is a set of certificates the attorney office can review as a specialised book, and it is the set in which the stock sits.
The same percentage is not equally informative inside each segment.
We therefore report the cut for the pooled book and again inside declared debt and inside assessed debt.
Nothing in the cutoff is a claim that working those certificates would recover the stock.

A firm has repeated enrolments when it holds two or more certificates in the segment.
We also record firms that hold at least one certificate in each segment.
The two counts answer different questions.
The first is about the book inside one door.
The second is about overlap.

\subsection{Comparisons}\label{sec:comparisons}

For a continuous outcome we use the Mann--Whitney test and the rank-biserial correlation \citep{mann1947,kerby2014}.
The correlation is $2U/(n_1 n_2) - 1$, where $U$ is the Mann--Whitney statistic for the assessed sample.
A positive value means that a random assessed unit tends to exceed a random declared unit.
We compute it for the certificate and for the firm.
The certificate-level statistic describes the portfolio.
Certificates of the same firm are not independent, so the firm-level correlation is the one we use when the claim is about debtors.
The samples are large.
Every $p$-value from these tests is zero at floating-point precision.
We report the effect size.

For a categorical attribute we use a chi-square test on firm counts and Cram\'er's $V$ \citep{cramer1946},
\begin{equation}
V = \sqrt{\chi^2 / \big(n \cdot (\min(r, c) - 1)\big)},
\label{eq:v}
\end{equation}
with $r$ and $c$ the dimensions of the table.
Each firm is placed in the segment that holds more of its stock, so a firm is one row.
The repeated-enrolment test is the exception: a firm is counted once in each segment in which it has a certificate, because the question is the shape of each book.
Levels that would leave a cell below five firms are pooled before the test.
$V$ near 0.1 is a small association and $V$ near 0.3 is a moderate one.
We use those conventional markers only as a reading aid.
